\section{Reduction on unions of members fails: the spine family}\label{app:spine}\label{sec:reduction}
Theorem~\ref{thm:three} applies Rule~6.2 of~\cite{BFK18} to components of $I(H)$. This appendix shows that the natural generalisation to pair-closed unions of \emph{members} (Section~\ref{sec:K}) cannot prove the reduction lemma (Open Problem~\ref{op:ML}), already for $w=1$ and $\ell=1$. Soundness of the rule is cited from~\cite{BFK18}; the proposition itself is proved here, and does not use the rule. Throughout, $\Gamma'$, $\mathrm{BIG}$, $B$, $Z$ are as in Section~\ref{sec:nl}; a set $\mathcal F$ of members is pair-closed, reducible and irreducible as in Section~\ref{sec:K}.

\begin{proposition}[The spine family $\mathrm{Sp}(F)$]\label{prop:five}
Let $F$ be a cubic graph with an orientation of in-degree at most $2$ at every vertex. Let $F'$ be $F$ plus a new vertex $*$ and an edge $e_0=(t\to *)$, $t\in V(F)$. Order the edges of $F'$ as $e_0,e_1,\dots,e_N$ so that $e_0$ comes first, \emph{consecutive edges have different heads and $\mathrm{head}(e_1)\ne\mathrm{tail}(e_0)=t$} (possible, since every head occurs at most twice and $t$ has in-degree at most $2$). Write $h(i)$ and $\mathrm{t}(i)$ for the head and tail of $e_i$. The instance $\mathrm{Sp}(F)$ has vertices $a_i,b_i,c_i,u_i$ ($0\le i\le N$), $z$ and $y$; for $v\in V(F')$ the sets $\mathcal K_v^+=\{a_i:\mathrm t(i)=v\}\cup\{b_i,c_i,u_i:h(i)=v\}$ (so $\mathcal K_*^+=\{b_0,c_0,u_0\}$). Edges: all pairs inside each $\mathcal K_v^+$; the cross edges $a_ib_i$; the spine edges $b_iu_{i+1}$ ($i<N$) and $b_Nz$; and $yu_0$, $yc_0$ (and no other edge at $y$). Packing: $P_i=(a_i,b_i,c_i)$ and $X_i=(u_i,b_i,u_{i+1})$, both with centre $b_i$ ($u_{N+1}:=z$); $\ell=1$. Then:
\begin{enumerate}
\item[(S1)] the members are induced $P_3$'s with pairwise disjoint pair sets, $H$ is pair-maximal, $I(H)$ is a caterpillar (spine $X_0,\dots,X_N$, leaf $P_i$ at $X_i$), so $w=1$, $c=2$, $D=4$;
\item[(S2)] $(G,H,1)$ is a YES instance and $\OPT=|H|+1$;
\item[(S3)] all hard degrees are at most $7\le B=11$, so $\mathrm{BIG}=\emptyset$;
\item[(S4)] $F$ is a minor of $\Gamma'$, so $\tw(\Gamma')\ge\tw(F)$, which is $\Omega(|V(F)|)$ for cubic expanders~\cite{GM09} (cited);
\item[(S5)] $I(H)$ is connected and tight: $\tau=|H|$;
\item[(S6)] every nonempty pair-closed set $\mathcal F$ of members (connected or not) is irreducible.
\end{enumerate}
\end{proposition}
\begin{proof}
(S1) Every induced $P_3$ of $G$ uses a cross edge or a spine edge or a $y$-edge, because two edges inside cliques at a common vertex lie in one clique. $P_i$ is induced: $a_ib_i$ is a cross edge, $b_ic_i$ lies in $\mathcal K^+_{h(i)}$, $a_ic_i\notin E$ because $\mathrm t(i)\ne h(i)$. $X_i$ is induced: $u_ib_i$ lies in $\mathcal K^+_{h(i)}$, $b_iu_{i+1}$ is a spine edge, and $u_iu_{i+1}\notin E$ since consecutive heads differ ($z$ is adjacent to $b_N$ only). The pair sets $\{a_ib_i,b_ic_i,a_ic_i\}$, $\{u_ib_i,b_iu_{i+1},u_iu_{i+1}\}$ are pairwise disjoint. Pair-maximality: no two hard edges form an induced $P_3$. The soft edges are $a_ib_i$, $b_ic_i$, $u_ib_i$, $b_iu_{i+1}$, so the hard edges are the other clique edges and $yu_0$, $yc_0$. Two hard clique edges at a vertex lie in one clique, so their other ends are adjacent; the hard edges at $y$ go to $u_0,c_0$, which are adjacent; and a hard edge at $u_0$ or $c_0$ inside $\mathcal K_*^+=\{b_0,c_0,u_0\}$ can only be $u_0c_0$, whose ends are both adjacent to $y$.

(S2) The partition $\{\mathcal K_v^+\}$ together with $\{z\}$, with $y$ put into $\mathcal K_*^+$, deletes exactly the $2(N+1)=|H|$ cross and spine edges (all soft) and inserts the single pair $yb_0$; its cost is $|H|+1$. For $\OPT\ge|H|+1$ see (S6) below.

(S3) A vertex of $\mathcal K_v^+$ has at most $|\mathcal K_v^+|-1=2+2\,\mathrm{in}(v)\le6$ hard neighbours inside its clique, and $u_0,c_0$ have one more hard neighbour, $y$ (which has hard degree $2$).

(S4) Contract every $\mathcal K_v^+$ (a clique). The cross edges $a_ib_i$ join $\mathcal K^+_{\mathrm t(i)}$ and $\mathcal K^+_{h(i)}$, so the quotient contains $F$.

(S5) In the partition of (S2), the pair $a_ib_i$ is deleted, $b_ic_i$ kept, $a_ic_i$ a non-edge in different clusters; $u_ib_i$ kept, $b_iu_{i+1}$ deleted, $u_iu_{i+1}$ a non-edge in different clusters. So its restriction to $V^-=V\setminus\{y\}$ edits exactly one pair per member: $\tau\le|H|$, hence $\tau=|H|$.

(S6) For pair-closed $\mathcal F$ the restriction of the solution of (S2) to $W_{\mathcal F}$ has at most $|\mathcal F|$ edits (no member outside $\mathcal F$ has a pair inside $W_{\mathcal F}$, and $yb_0$ leaves $W_{\mathcal F}$), so $\tau(G[W_{\mathcal F}])=|\mathcal F|$ and the optimal $S_{\mathcal F}$ are the \emph{tight} ones: exactly one edited pair per member of $\mathcal F$ and no hard pair is edited. Suppose such an $S_{\mathcal F}$ satisfied the boundary condition. We use repeatedly: two vertices $p,q\in W_{\mathcal F}$ joined by a hard non-edge are non-adjacent in $G[W_{\mathcal F}]\sd S_{\mathcal F}$, and a hard edge is kept.

\emph{Case 1: $\mathcal F$ contains a spine member; let $X_i$ be the one of smallest index.}

\emph{$i\ge1$.} Then $X_{i-1}\notin\mathcal F$, and $P_{i-1}\notin\mathcal F$: otherwise $b_{i-1},u_i\in W_{\mathcal F}$ and the pair $b_{i-1}u_i$ of the outside member $X_{i-1}$ would lie inside $W_{\mathcal F}$. So $a_{i-1},b_{i-1},u_{i-1}\notin W_{\mathcal F}$ and $N(b_{i-1})\cap W_{\mathcal F}=\{u_i\}\cup Y$ with $Y=\mathcal K^+_{h(i-1)}\cap W_{\mathcal F}$. For $m\in Y$ the pair $u_im$ is a non-edge ($h(i)\ne h(i-1)$) and is soft only if $m=u_{i+1}$. If $Y$ contains some $m\ne u_{i+1}$, then $u_i,m$ are two non-adjacent vertices of $N(b_{i-1})\cap W_{\mathcal F}$: the boundary condition fails. If $Y=\{u_{i+1}\}$, the set $\{u_i,u_{i+1}\}$ must be a cluster, so $X_i$ inserts $u_iu_{i+1}$ and keeps $u_ib_i$, $b_iu_{i+1}$; then $b_i$ lies in this cluster, but $b_i\notin N(b_{i-1})$: contradiction. So $Y=\emptyset$ and $\{u_i\}$ is a cluster, i.e.\ $X_i$ deletes $u_ib_i$. If $P_i\notin\mathcal F$, the vertex $c_i\notin W_{\mathcal F}$ is adjacent to $u_i$ and $b_i$, which lie in different clusters: the boundary condition fails at $c_i$. If $P_i\in\mathcal F$, the hard edge $u_ic_i$ is kept, contradicting that $u_i$ is a singleton.

\emph{$i=0$.} The vertex $y\notin W_{\mathcal F}$ has $N(y)\cap W_{\mathcal F}=\{u_0\}$ if $P_0\notin\mathcal F$ and $\{u_0,c_0\}$ if $P_0\in\mathcal F$. First case: $\{u_0\}$ is a cluster, so $X_0$ deletes $u_0b_0$, and $c_0\notin W_{\mathcal F}$, adjacent to $u_0,b_0\in W_{\mathcal F}$ in different clusters, violates the condition. Second case: $\{u_0,c_0\}$ must be a cluster, so $b_0$ is separated from both: $X_0$ deletes $u_0b_0$ and $P_0$ deletes $b_0c_0$. Then $a_0b_0$ and $b_0u_1$ are kept, so $a_0,u_1$ are adjacent to $b_0$. Now $a_0\in \mathcal K_t^+$ and $u_1\in \mathcal K^+_{h(1)}$ with $h(1)\ne t$ \emph{by the ordering condition}, so $a_0u_1$ is a hard non-edge (the soft partners of $a_0$ are $b_0,c_0$), and $a_0b_0u_1$ is an induced $P_3$ of $G[W_{\mathcal F}]\sd S_{\mathcal F}$: contradiction.

\emph{Case 2: $\mathcal F$ contains only leaves.} Take $P_i\in\mathcal F$; $u_i,u_{i+1}\notin W_{\mathcal F}$ (for $i=N$ read $z$ for $u_{i+1}$). If $i=N$ then $N(z)\cap W_{\mathcal F}=\{b_N\}$ must be a cluster, i.e.\ $b_N$ is a singleton; but each of the three single edits of $P_N$ (delete $a_Nb_N$, delete $b_Nc_N$, insert $a_Nc_N$) keeps $b_N$ with $a_N$ or $c_N$. So let $i<N$. Then $N(u_{i+1})\cap W_{\mathcal F}=\{b_i\}\cup Y$ with $Y=\mathcal K^+_{h(i+1)}\cap W_{\mathcal F}$. For $m\in Y$ the pair $b_im$ is a non-edge ($h(i+1)\ne h(i)$) unless $m=a_i$ (possible when $\mathrm t(i)=h(i+1)$), and it is soft only if $m=a_i$ (the soft partners of $b_i$ are $a_i,c_i,u_i,u_{i+1}$). \emph{(2a)} If $Y$ contains some $m\ne a_i$, then $b_i,m$ are non-adjacent in $N(u_{i+1})\cap W_{\mathcal F}$: failure. \emph{(2b)} If $Y=\emptyset$, then $\{b_i\}$ is a cluster, impossible as above. \emph{(2c)} If $Y=\{a_i\}$, then $P_{i+1}\notin\mathcal F$ and $b_{i+1}\in \mathcal K^+_{h(i+1)}$ is outside $W_{\mathcal F}$ with $N(b_{i+1})\cap W_{\mathcal F}=\{a_i\}$; this must be a cluster, so $a_i$ is a singleton and $a_ib_i$ is deleted. But then $\{a_i,b_i\}=N(u_{i+1})\cap W_{\mathcal F}$ is not a cluster ($a_ib_i$ is not an edge of $G\sd S_{\mathcal F}$): contradiction.

So $\mathcal F$ is irreducible. For the lower bound of (S2): if $\OPT=|H|$, an optimum $S$ has $|S|=\tau$ for the whole component $\mathcal C=I(H)$, so $S=S_{\mathcal C}$ is tight and optimal on $V^-=V\setminus\{y\}$; by (S6) it violates the boundary condition at $y$, and Lemma~\ref{lem:D} gives an edited pair $\{y,u\}$, which contradicts $S=S_{\mathcal C}$. So $\OPT\ge|H|+1$.
\end{proof}
\begin{remark}[Finite checks]\label{rem:fivecov}
Claims (S1)--(S5) and $\OPT=|H|+1$ were verified by computer for $F$ the triangle ($n=18$, $\OPT$ by an exact subset dynamic program) and for $K_4$ ($n=30$). Claim (S6) was checked, in agreement with the proof, for the triangle on all $90$ pair-closed subfamilies with $|\mathcal F|\le6$ and on the whole component, and for $K_4$ on all $632$ pair-closed subfamilies with $|\mathcal F|\le4$ for $8$ random admissible orderings. The proof above is the justification.
\end{remark}

\paragraph{Consequence: an obstruction, not a refutation of the reduction lemma.}
No reduction obtained by iterating Rule~6.2 on pair-closed unions of members (the form used for components in Section~\ref{sec:K}) can prove Open Problem~\ref{op:ML} by itself: on $\mathrm{Sp}(F)$ it removes nothing, although $\ell=1$, $w=1$, $\mathrm{BIG}=\emptyset$ and $\tw(\Gamma')=\Omega(|V(F)|)$. In particular the counting argument of Lemma~\ref{lem:count} does not extend from components to members: all $2(N+1)$ members are irreducible although $\ell=1$. The mechanism is Lemma~\ref{lem:D}: along the spine each member is made irreducible by an edit at a shared vertex of its neighbour, down to the trigger $y$. Other generalisations of the rule (to vertex sets that are not unions of members) are not excluded.

